%% kapa: exam=1 round=29 year=2018 period=1 form=A subject=real_estate subject_ko=부동산학원론 seq=12 number=92
%% source: Q-Net 감정평가사 기출문제 (cst00309) artlSeq=5202808 file=extracted/1008/5202808/2018년 제29회 감정평가사 1차시험 기출문제/2018년 제29회 감정평가사 1차시험 기출문제/1교시 A형.pdf
%% conversion: kordoc 4.24.51 -> tools/gen1.py (auto); review=proofread-sonnet; flags=none
\begin{question}{92}
다음과 같은 상황이 주어졌을 때 총투자수익률(ROI: return on investment)과 부채감당률(DCR: debt coverage ratio)은? (단, 총투자기간은 1년, 수치는 소수점 이하 둘째자리에서 반올림함)
\begin{framedtext}
○ 총투자액: 12억원(현금) \par
○ 순영업소득(순수익): 1억원/연 \par
○ 세전현금흐름(세전현금수지): 7천만원/연 \par
○ 저당지불액(부채서비스액): 5천만원/연 \par
○ 공실 및 불량부채(공실손실상당액 및 대손충당금): 5천만원/연 \par
○ 사업소득세: 1천만원/연 \par
○ 유효총소득승수: 4 \par
○ 지분투자액: 5억원 \par
\end{framedtext}
\begin{choices}
  \item 총투자수익률: 29.2\%, 부채감당률: 2.0
  \item 총투자수익률: 29.2\%, 부채감당률: 2.4
  \item 총투자수익률: 33.3\%, 부채감당률: 2.0
  \item 총투자수익률: 33.3\%, 부채감당률: 2.4
  \item 총투자수익률: 37.2\%, 부채감당률: 3.0
\end{choices}
\end{question}
\answerall
\answernote{정답표에 1,2,3,4,5 모두 정답으로 표기(전원 정답)}
